% Source slide 18
\section[4.2 Null and Column Spaces]{4.2 Null Spaces, Column Spaces, and Linear Transformations}
\chaptercard{4}{Vector Spaces}{4.2}{NULL SPACES, COLUMN SPACES, AND LINEAR TRANSFORMATIONS}

% Source slide 19
\begin{frame}[shrink=15]{NULL SPACE OF A MATRIX}
\begin{itemize}
\item \textbf{Definition:} The \textcolor{sourceblue}{\textbf{null space}} of an $m\times n$ matrix $A$, written as $\Nul A$, is the set of all solutions of the homogeneous equation $A\vect x=0$. In set notation,
\[
\Nul A=\{\vect x:\vect x\text{ is in }\R^n\text{ and }A\vect x=0\}.
\]
\item \sourcehead{Theorem 2:} The null space of an $m\times n$ matrix $A$ is a subspace of $\R^n$. Equivalently, the set of all solutions to a system $A\vect x=0$ of $m$ homogeneous linear equations in $n$ unknowns is a subspace of $\R^n$.
\item \textbf{Proof:} $\Nul A$ is a subset of $\R^n$ because $A$ has $n$ columns.
\item We need to show that $\Nul A$ satisfies the three properties of a subspace.
\end{itemize}
\end{frame}

% Source slide 20
\begin{frame}[shrink=13]{NULL SPACE OF A MATRIX}
\begin{itemize}
\item $0$ is in $\Nul A$.
\item Next, let $\vect u$ and $\vect v$ represent any two vectors in $\Nul A$.
\item Then
\[
A\vect u=0\qquad\text{and}\qquad A\vect v=0.
\]
\item To show that $\vect u+\vect v$ is in $\Nul A$, we must show that $A(\vect u+\vect v)=0$.
\item Using a property of matrix multiplication, compute
\[
A(\vect u+\vect v)=A\vect u+A\vect v=0+0=0.
\]
\item Thus $\vect u+\vect v$ is in $\Nul A$, and $\Nul A$ is closed under vector addition.
\end{itemize}
\end{frame}

% Source slide 21
\begin{frame}[shrink=10]{NULL SPACE OF A MATRIX}
\begin{itemize}
\item Finally, if $c$ is any scalar, then
\[
A(c\vect u)=c(A\vect u)=c(0)=0,
\]
which shows that $c\vect u$ is in $\Nul A$.
\item Thus $\Nul A$ is a subspace of $\R^n$.
\item \textbf{An Explicit Description of $\Nul A$}
\item There is no obvious relation between vectors in $\Nul A$ and the entries in $A$.
\item We say that $\Nul A$ is defined \emph{implicitly}, because it is defined by a condition that must be checked.
\end{itemize}
\end{frame}

% Source slide 22
\begin{frame}[shrink=10]{NULL SPACE OF A MATRIX}
\begin{itemize}
\item No explicit list or description of the elements in $\Nul A$ is given.
\item \emph{Solving} the equation $A\vect x=0$ amounts to producing an \emph{explicit} description of $\Nul A$.
\item \textbf{Example 3:} Find a spanning set for the null space of the matrix
\[
A=\begin{bmatrix}
-3&6&-1&1&-7\\
1&-2&2&3&-1\\
2&-4&5&8&-4
\end{bmatrix}.
\]
\end{itemize}
\end{frame}

% Source slide 23
\begin{frame}[shrink=15]{NULL SPACE OF A MATRIX}
\begin{itemize}
\item \textbf{Solution:} The first step is to find the general solution of $A\vect x=0$ in terms of free variables.
\item Row reduce the augmented matrix $[A\ \ 0]$ to reduce echelon form in order to write the basic variables in terms of the free variables:
\end{itemize}
\[
\begin{bmatrix}
1&-2&0&-1&3&0\\
0&0&1&2&-2&0\\
0&0&0&0&0&0
\end{bmatrix},
\qquad
\begin{aligned}
x_1-2x_2-x_4+3x_5&=0,\\
x_3+2x_4-2x_5&=0,\\
0&=0.
\end{aligned}
\]
\end{frame}

% Source slide 24
\begin{frame}[shrink=18]{NULL SPACE OF A MATRIX}
\begin{itemize}
\item The general solution is $x_1=2x_2+x_4-3x_5$, $x_3=-2x_4+2x_5$, with $x_2,x_4$, and $x_5$ free.
\item Next, decompose the vector giving the general solution into a linear combination of \emph{vectors where the weights are the free variables}. That is,
\end{itemize}
\[
\begin{bmatrix}x_1\\x_2\\x_3\\x_4\\x_5\end{bmatrix}
=\begin{bmatrix}2x_2+x_4-3x_5\\x_2\\-2x_4+2x_5\\x_4\\x_5\end{bmatrix}
=x_2\underbrace{\begin{bmatrix}2\\1\\0\\0\\0\end{bmatrix}}_{\vect u}
+x_4\underbrace{\begin{bmatrix}1\\0\\-2\\1\\0\end{bmatrix}}_{\vect v}
+x_5\underbrace{\begin{bmatrix}-3\\0\\2\\0\\1\end{bmatrix}}_{\vect w}.
\]
\end{frame}

% Source slide 25
\begin{frame}[shrink=15]{NULL SPACE OF A MATRIX}
\[
=x_2\vect u+x_4\vect v+x_5\vect w \tag{3}
\]
\begin{itemize}
\item Every linear combination of $\vect u$, $\vect v$, and $\vect w$ is an element of $\Nul A$.
\item Thus $\{\vect u,\vect v,\vect w\}$ is a spanning set for $\Nul A$.
\end{itemize}
\begin{enumerate}
\item The spanning set produced by the method in Example (3) is automatically linearly independent because the free variables are the weights on the spanning vectors.
\item When $\Nul A$ contains nonzero vectors, the number of vectors in the spanning set for $\Nul A$ equals the number of free variables in the equation $A\vect x=0$.
\end{enumerate}
\end{frame}

% Source slide 26
\begin{frame}[shrink=13]{COLUMN SPACE OF A MATRIX}
\begin{itemize}
\item \textbf{Definition:} The column space of an $m\times n$ matrix $A$, written as $\Col A$, is the set of all linear combinations of the columns of $A$. If $A=[\vect a_1\ \cdots\ \vect a_n]$, then
\[
\Col A=\Span\{\vect a_1,\ldots,\vect a_n\}.
\]
\item \sourcehead{Theorem 3:} The column space of an $m\times n$ matrix $A$ is a subspace of $\R^m$.
\item A typical vector in $\Col A$ can be written as $A\vect x$ for some $\vect x$ because the notation $A\vect x$ stands for a linear combination of the columns of $A$. That is,
\[
\Col A=\{\vect b:\vect b=A\vect x\text{ for some }\vect x\text{ in }\R^n\}.
\]
\end{itemize}
\end{frame}

% Source slide 27
\begin{frame}[shrink=13]{COLUMN SPACE OF A MATRIX}
\begin{itemize}
\item The column space of an $m\times n$ matrix $A$ is all of $\R^m$ if and only if the equation $A\vect x=\vect b$ has a solution for each $\vect b$ in $\R^m$.
\item \textbf{Example 7:} Let
\[
A=\begin{bmatrix}2&4&-2&1\\-2&-5&7&3\\3&7&-8&6\end{bmatrix},
\qquad
\vect u=\begin{bmatrix}3\\-2\\-1\\0\end{bmatrix},
\qquad
\vect v=\begin{bmatrix}3\\-1\\3\end{bmatrix}.
\]
\end{itemize}
\end{frame}

% Source slide 28
\begin{frame}[shrink=14]{COLUMN SPACE OF A MATRIX}
\begin{enumerate}
\item[(a.)] Determine if $\vect u$ is in $\Nul A$. Could $\vect u$ be in $\Col A$?
\item[(b.)] Determine if $\vect v$ is in $\Col A$. Could $\vect v$ be in $\Nul A$?
\end{enumerate}
\begin{itemize}
\item \textbf{Solution:}
\end{itemize}
\begin{enumerate}
\item[(a.)] An explicit description of $\Nul A$ is not needed here. Simply compute the product $A\vect u$.
\end{enumerate}
\[
A\vect u=
\begin{bmatrix}2&4&-2&1\\-2&-5&7&3\\3&7&-8&6\end{bmatrix}
\begin{bmatrix}3\\-2\\-1\\0\end{bmatrix}
=\begin{bmatrix}0\\-3\\3\end{bmatrix}
\ne\begin{bmatrix}0\\0\\0\end{bmatrix}.
\]
\end{frame}

% Source slide 29
\begin{frame}[shrink=14]{COLUMN SPACE OF A MATRIX}
\begin{itemize}
\item $\vect u$ is \emph{not} a solution of $A\vect x=0$, so $\vect u$ is not in $\Nul A$.
\item Also, with four entries, $\vect u$ could not possibly be in $\Col A$, since $\Col A$ is a subspace of $\R^3$.
\end{itemize}
\begin{enumerate}
\item[(b.)] Reduce $[A\ \ \vect v]$ to an echelon form.
\end{enumerate}
\[
[A\ \ \vect v]=
\begin{bmatrix}2&4&-2&1&3\\-2&-5&7&3&-1\\3&7&-8&6&3\end{bmatrix}
\sim
\begin{bmatrix}2&4&-2&1&3\\0&1&-5&-4&-2\\0&0&0&17&1\end{bmatrix}.
\]
\begin{itemize}
\item The equation $A\vect x=\vect v$ is consistent, so $\vect v$ is in $\Col A$.
\end{itemize}
\end{frame}

% Source slide 30
\begin{frame}[shrink=13]{CONTRAST BETWEEN $\Nul A$ AND $\Col A$ FOR AN $m\times n$ MATRIX $A$}
\renewcommand{\arraystretch}{1.35}
\begin{tabular}{|>{\raggedright\arraybackslash}p{.45\linewidth}|>{\raggedright\arraybackslash}p{.45\linewidth}|}
\hline
\centering $\Nul A$ & \centering $\Col A$\tabularnewline
\hline
1. $\Nul A$ is a subspace of $\R^n$. & 1. $\Col A$ is a subspace of $\R^m$.\\
\hline
2. $\Nul A$ is implicitly defined; \emph{i.e.}, you are given only a condition $(A\vect x=0)$ that vectors in $\Nul A$ must satisfy. & 2. $\Col A$ is explicitly defined; \emph{i.e.}, you are told how to build vectors in $\Col A$.\\
\hline
\end{tabular}
\end{frame}

% Source slide 31
\begin{frame}[shrink=13]{CONTRAST BETWEEN $\Nul A$ AND $\Col A$ FOR AN $m\times n$ MATRIX $A$}
\renewcommand{\arraystretch}{1.35}
\begin{tabular}{|>{\raggedright\arraybackslash}p{.45\linewidth}|>{\raggedright\arraybackslash}p{.45\linewidth}|}
\hline
3. It takes time to find vectors in $\Nul A$. Row operations on $[A\ \ 0]$ are required. & 3. It is easy to find vectors in $\Col A$. The columns of a are displayed; others are formed from them.\\
\hline
4. There is no obvious relation between $\Nul A$ and the entries in $A$. & 4. There is an obvious relation between $\Col A$ and the entries in $A$, since each column of $A$ is in $\Col A$.\\
\hline
\end{tabular}
\end{frame}

% Source slide 32
\begin{frame}[shrink=13]{CONTRAST BETWEEN $\Nul A$ AND $\Col A$ FOR AN $m\times n$ MATRIX $A$}
\renewcommand{\arraystretch}{1.35}
\begin{tabular}{|>{\raggedright\arraybackslash}p{.45\linewidth}|>{\raggedright\arraybackslash}p{.45\linewidth}|}
\hline
5. A typical vector $\vect v$ in $\Nul A$ has the property that $A\vect v=0$. & 5. A typical vector $\vect v$ in $\Col A$ has the property that the equation $A\vect x=\vect v$ is consistent.\\
\hline
6. Given a specific vector $\vect v$, it is easy to tell if $\vect v$ is in $\Nul A$. Just compare $A\vect v$. & 6. Given a specific vector $\vect v$, it may take time to tell if $\vect v$ is in $\Col A$. Row operations on $[A\ \ \vect v]$ are required.\\
\hline
\end{tabular}
\end{frame}

% Source slide 33
\begin{frame}[shrink=13]{CONTRAST BETWEEN $\Nul A$ AND $\Col A$ FOR AN $m\times n$ MATRIX $A$}
\renewcommand{\arraystretch}{1.35}
\begin{tabular}{|>{\raggedright\arraybackslash}p{.45\linewidth}|>{\raggedright\arraybackslash}p{.45\linewidth}|}
\hline
7. $\Nul A=\{0\}$ if and only if the equation $A\vect x=0$ has only the trivial solution. & 7. $\Col A=\R^m$ if and only if the equation $A\vect x=\vect b$ has a solution for every $\vect b$ in $\R^m$.\\
\hline
8. $\Nul A=\{0\}$ if and only if the linear transformation $\vect x\mapsto A\vect x$ is one-to-one. & 8. $\Col A=\R^m$ if and only if the linear transformation $\vect x\mapsto A\vect x$ maps $\R^n$ onto $\R^m$.\\
\hline
\end{tabular}
\end{frame}

% Source slide 34
\begin{frame}{KERNEL AND RANGE OF A LINEAR TRANSFORMATION}
\begin{itemize}
\item The \textbf{kernel} (or \textbf{null space}) of such a $T$ is the set of all $\vect u$ in $V$ such that $T(\vect u)=0$ (the zero vector in $W$).
\item The \textbf{range} of $T$ is the set of all vectors in $W$ of the form $T(\vect x)$ for some $\vect x$ in $V$.
\item The kernel of $T$ is a subspace of $V$.
\item The range of $T$ is a subspace of $W$.
\end{itemize}
\end{frame}
